\documentclass{article}
\usepackage[T1]{fontenc}
\usepackage{lmodern}
\usepackage{graphicx}
\usepackage{amsmath,amssymb}
\usepackage[a4paper,margin=2cm]{geometry}
\usepackage[absolute,overlay]{textpos}
\setlength{\TPHorizModule}{1mm}
\setlength{\TPVertModule}{1mm}
\usepackage[indonesian]{babel}
\usepackage{hyperref}
\setlength{\emergencystretch}{2em}
% unit: Halaman 63

\begin{document}

% === Halaman 63 (python/native) ===

• Biasanya, plainteks dimanipulasi dulu dengan kunci (umumnya di-XOR-kan 

dengan kunci eksternal, atau dengan kunci putaran pertama), lalu hasilnya 

ditransformasikan dengan fungsi f berulang kali.




          𝐶0 = P   K0 

                       𝐶𝑖 = 𝑓(𝐶𝑖−1, 𝐾𝑖) 

, 𝑖= 1, 2, … , 𝑟−1



𝐶= 𝐶𝑟 = 𝐶𝑟−1  𝐾𝑟

• Semakin banyak jumlah putaran, efek diffusi semakin bertambah, namun jumlah 

putaran yang besar dapat mempuat komputasi menjadi tidak sangkil (efisien).

• Jumlah putaran harus dipertimbangkan untuk mempertemukan kriteria 

keamanan dan efisiensi. 


63

\end{document}
